\scanpage{124}
\mspage{111}
et de $=$ pour $B,X,Y$. Plus précisément,
la catégorie $M(I^\circ)$ s'interprète comme
celle des triples $X=(X',X_0,\alpha_X)$, où
$X'\in M(I'{}^\circ)$, $X_0\in M$, et $\alpha_X\colon X'_{I'}\longrightarrow X_0$.

À ce compte, la donnée du diagramme
(6.3.2) équivaut à celle d'un diagramme \inserted{carré}
\unsure{similaire|semblable}

\diagram{\((6.3.3)\)\qquad
\begin{tikzcd}[ampersand replacement=\&]
 A'\arrow[r,"{\alpha'}"]\arrow[d,"{i'}"'] \& X'\arrow[d,"{p'}"]\\
 B'\arrow[r,"{\beta'}"]\arrow[ur,dotted,"{h'}"] \& Y'
\end{tikzcd}}

\uline{dans $M(I'{}^\circ)$}, \unsure{d'où|d'un} \unsure{autre|un} diagramme
dans $M$

\diagram{\((6.3.4)\)\qquad
\begin{tikzcd}[ampersand replacement=\&]
 A'_{I'}\arrow[r,"{\alpha'_{I'}}"]\arrow[d,"{i'_{I'}}"'] \& X'_{I'}\arrow[d,"{p'_{I'}}"]\\
 B'_{I'}\arrow[r,"{\beta'_{I'}}"]\arrow[ur,dotted,"{h'_{I'}}"] \& Y'_{I'},
\end{tikzcd}}

d'un \inserted{autre} diagramme carré \uline{dans $M$}

\diagram{\((6.3.5)\)\qquad
\begin{tikzcd}[ampersand replacement=\&]
 A_0\arrow[r,"{\alpha_0}"]\arrow[d,"{i_0}"'] \& X_0\arrow[d,"{p_0}"]\\
 B_0\arrow[r,"{\beta_0}"] \& Y_0
\end{tikzcd}}

et enfin d'un hom. \inserted{$\lambda=(\lambda_A,\lambda_B,\lambda_X,\lambda_Y))$} du carré (6.3.4)
dans le carré (6.3.5.) : \deleted{\unsure{la}}
% archive sheet number: 123

